\documentclass[11pt,a4paper]{article}

\usepackage{fontspec}
\setmainfont{DejaVu Serif}[
  BoldFont={DejaVu Serif Bold},
  ItalicFont={DejaVu Serif},
  ItalicFeatures={FakeSlant=0.18},
  BoldItalicFont={DejaVu Serif Bold},
  BoldItalicFeatures={FakeSlant=0.18}
]
\setsansfont{DejaVu Sans}
\usepackage{polyglossia}
\setmainlanguage{english}
\usepackage[a4paper,margin=21mm]{geometry}
\usepackage{microtype}
\usepackage{setspace}
\setstretch{1.035}
\usepackage{amsmath,amssymb,amsthm,mathtools,bm,mathrsfs}
\usepackage{booktabs,longtable,array,tabularx}
\usepackage{enumitem}
\usepackage{xcolor}
\usepackage{tikz}
\usetikzlibrary{arrows.meta,positioning,calc}
\usepackage[unicode,hidelinks]{hyperref}
\usepackage[nameinlink,noabbrev]{cleveref}
\emergencystretch=3em
\allowdisplaybreaks

\hypersetup{
  pdftitle={Distance and Scale Interface under Observation Maps},
  pdfauthor={Stassis Research Program},
  pdfsubject={GO Core v0.2 reference module for observed pseudometrics, quotient metrics, shells, resolutions, and unit-distance graphs},
  pdfkeywords={metric, pseudometric, observation map, quotient metric, scale, resolution, unit distance, Geometry of Observation}
}

\definecolor{obsblue}{RGB}{37,83,138}
\definecolor{obsred}{RGB}{150,45,45}
\definecolor{obsgreen}{RGB}{30,115,78}
\definecolor{lightblue}{RGB}{239,246,253}
\definecolor{lightred}{RGB}{253,242,242}
\definecolor{lightgreen}{RGB}{239,249,244}

\theoremstyle{definition}
\newtheorem{definition}{Definition}[section]
\newtheorem{model}[definition]{Model}
\theoremstyle{plain}
\newtheorem{theorem}[definition]{Theorem}
\newtheorem{proposition}[definition]{Proposition}
\newtheorem{lemma}[definition]{Lemma}
\newtheorem{corollary}[definition]{Corollary}
\theoremstyle{remark}
\newtheorem{remark}[definition]{Remark}
\newtheorem{warning}[definition]{Boundary}

\newcommand{\R}{\mathbb{R}}
\newcommand{\Z}{\mathbb{Z}}
\newcommand{\X}{\mathcal{X}}
\newcommand{\Y}{\mathcal{Y}}
\newcommand{\Ocal}{\mathcal{O}}
\newcommand{\Pcal}{\mathcal{P}}
\newcommand{\Bcal}{\mathcal{B}}
\newcommand{\dd}{\mathop{}\!d}
\newcommand{\dist}{\operatorname{dist}}
\newcommand{\diam}{\operatorname{diam}}
\newcommand{\Lip}{\operatorname{Lip}}
\newcommand{\im}{\operatorname{im}}
\newcommand{\card}{\operatorname{card}}
\newcommand{\abs}[1]{\left\lvert #1\right\rvert}
\newcommand{\norm}[1]{\left\lVert #1\right\rVert}
\newcommand{\set}[1]{\left\{#1\right\}}
\newcommand{\status}[1]{\textsf{\footnotesize #1}}

\newcolumntype{L}[1]{>{\raggedright\arraybackslash}p{#1}}
\newcolumntype{Y}{>{\raggedright\arraybackslash}X}

\title{\bfseries Distance and Scale Interface under Observation Maps\\
\large Global pseudometrics, quotient geometry, matched resolutions, and calibrated shells}
\author{Stassis Research Program}
\date{Strict GO Core reference revision v0.2 --- 2026}

\begin{document}
\maketitle

\begin{center}
\fcolorbox{obsblue}{lightblue}{%
\parbox{0.91\linewidth}{\small
\textbf{Status.} This document replaces \emph{Unit Distances under
Observation Maps v0.1} as the distance-scale reference module for
GO Core v0.2. A numerical value ``1'' is never compared with a physical
distance without a declared reference radius. The global observed
pseudometric is separated from the local pullback tensor, all resolution
maps are typed, and exact-shell statements are separated from tolerant or
noisy decisions.}}
\end{center}

\begin{abstract}
Let $O:(X,d_X)\to(Y,d_Y)$ be an observation map between metric spaces.
The globally observed distance on hidden states is the pseudometric
$d_O(x,x')=d_Y(Ox,Ox')$. Its zero classes are the observation fibers, and
the metric quotient is canonically isometric to the image $O(X)$. This
global construction is distinct from the local pullback quadratic form
$O^*g_Y$; the corresponding path pseudometric can exceed the ambient
image distance unless a minimizing image geodesic is liftable. A
distance shell is defined relative to a positive, typed radius $r_*$,
while finite resolution and measurement error require a tolerance
$\tau$ and an explicit decision rule. Under coherent changes of unit,
$d_Y,r_*,\tau$, and the output resolution rescale together, leaving all
normalized predicates invariant. Lipschitz observations propagate input
resolution to output resolution and compose multiplicatively. The
orthogonal-projection cylinder, fiber multiplicities in observed
unit-distance graphs, and periodic half-step aliasing are derived as
applications. The module is an interface built from standard metric and
differential geometry; it does not solve the classical Erdos unit-distance
problem.
\end{abstract}

\tableofcontents

\section{Scope, objects, and claim boundary}

The purpose of this module is to make every use of ``distance'' and
``scale'' in Geometry of Observation type-checkable. It supplies a common
interface for metric entropy, ruler traces, mechanics, and later
observation modules. It does not impose a physical ontology on the hidden
space.

\begin{warning}[Claim boundary]
An observation-induced pseudometric is a standard pullback construction.
Fiber multiplicity can create quadratically many hidden-state pairs, but
this is not a new lower bound for the classical planar unit-distance
function. A degenerate pullback tensor does not by itself imply a physical
extra dimension.
\end{warning}

\begin{definition}[Typed metric datum]
A typed metric datum is a tuple
\[
  \mathfrak M_X=(X,d_X,U_X),
\]
where $d_X:X\times X\to\R_{\geq0}U_X$ is a metric and $U_X$ names its
scale type. For a physical spatial metric, $U_X$ may be metres; for a
normalized abstract metric, $U_X=1$.
\end{definition}

The notation $\R_{\geq0}U_X$ is a bookkeeping device: it records that
$d_X$ is length-valued (or more generally scale-valued), not an
untyped real number. Ratios of two positive quantities with the same
scale type are dimensionless.

\begin{definition}[System and protocol passports]
The intrinsic descriptor and the observation protocol are distinct:
\[
 \xi_X=(X,d_X,\diam X,\text{intrinsic regularity data},\ldots),
\]
\[
 \lambda=(O,d_Y,r_*,\varepsilon_X,\varepsilon_Y,\tau,
          Q,\mathsf E,\mathsf U,\ldots).
\]
Here $r_*>0$ is a declared output-space reference radius,
$\varepsilon_X,\varepsilon_Y>0$ are input and output resolutions,
$\tau\geq0$ is a decision tolerance, $Q$ is a quantizer or partition,
$\mathsf E$ is an estimator or direct readout, and $\mathsf U$ is an
uncertainty model.
\end{definition}

An intrinsic diameter is not an instrument resolution. A shell radius is
not a regression reference length. The same physical dimension therefore
does not imply the same semantic role.

\section{The global observed pseudometric}

\begin{definition}[Observed pseudometric]
Let $(Y,d_Y)$ be a metric space and let $O:X\to Y$ be any map. Define
\[
  d_O(x,x'):=d_Y\bigl(O(x),O(x')\bigr).
\]
\end{definition}

\begin{theorem}[Pseudometric and separation criterion]
\label{thm:pseudometric}
The function $d_O$ is a pseudometric on $X$. It is a metric if and only if
$O$ is injective.
\end{theorem}

\begin{proof}
Nonnegativity and symmetry follow from $d_Y$. For $x,x',x''\in X$,
\[
d_O(x,x'')
=d_Y(Ox,Ox'')
\leq d_Y(Ox,Ox')+d_Y(Ox',Ox'')
=d_O(x,x')+d_O(x',x'').
\]
Thus the triangle inequality holds. Moreover,
$d_O(x,x')=0$ if and only if $O(x)=O(x')$. Hence separation is equivalent
to injectivity.
\end{proof}

\begin{definition}[Observation equivalence]
Write
\[
 x\sim_O x'
 \quad\Longleftrightarrow\quad
 O(x)=O(x').
\]
Let $X/{\sim_O}$ denote the set of observation fibers.
\end{definition}

\begin{theorem}[Canonical quotient metric]
\label{thm:quotient}
The formula
\[
 \bar d_O([x],[x']):=d_Y(Ox,Ox')
\]
is well-defined and is a metric on $X/{\sim_O}$. The map
\[
 \iota_O:X/{\sim_O}\longrightarrow O(X),\qquad [x]\longmapsto O(x),
\]
is an isometry onto the metric subspace $O(X)\subseteq Y$.
\end{theorem}

\begin{proof}
Changing representatives leaves both images unchanged, so the formula is
well-defined. The pseudometric axioms descend from
\cref{thm:pseudometric}; zero distance now implies equality of classes.
The map $\iota_O$ is a bijection and preserves the displayed distance.
\end{proof}

\begin{center}
\begin{tikzpicture}[
  node distance=12mm and 18mm,
  box/.style={draw,rounded corners,minimum width=29mm,minimum height=10mm,align=center},
  arr/.style={-{Latex[length=2.1mm]},thick}
]
\node[box,fill=lightblue] (x) {$X$\\hidden states};
\node[box,fill=lightred,right=of x] (q) {$X/{\sim_O}$\\observable classes};
\node[box,fill=lightgreen,right=of q] (y) {$O(X)\subseteq Y$\\metric image};
\draw[arr] (x) -- node[above] {quotient} (q);
\draw[arr] (q) -- node[above] {$\iota_O$ isometry} (y);
\draw[arr,obsred,bend right=24] (x.south) to node[below] {$O$} (y.south);
\end{tikzpicture}
\end{center}

\section{Local pullback form versus global distance}

Suppose now that $X$ and $Y$ are smooth manifolds, $Y$ has Riemannian
metric $g_Y$, and $O:X\to Y$ is smooth.

\begin{definition}[Pullback quadratic form]
\[
 g_O:=O^*g_Y,\qquad
 (g_O)_x(v,w)=
 (g_Y)_{O(x)}\bigl(\dd O_xv,\dd O_xw\bigr).
\]
\end{definition}

\begin{proposition}[Degeneracy directions]
$g_O$ is positive semidefinite and
\[
 \ker \dd O_x\subseteq\ker(g_O)_x.
\]
If $g_Y$ is positive definite, equality holds.
\end{proposition}

\begin{proof}
Positive semidefiniteness follows from
$g_O(v,v)=g_Y(\dd O\,v,\dd O\,v)\geq0$. If $\dd O\,v=0$, then
$g_O(v,w)=0$ for every $w$. Conversely, if $g_O(v,v)=0$, positive
definiteness of $g_Y$ gives $\dd O\,v=0$.
\end{proof}

For every piecewise $C^1$ curve $\gamma$ in $X$,
\[
 L_{g_O}(\gamma)=L_{g_Y}(O\circ\gamma).
\]
Define the extended path pseudometric
\[
 \rho_O(x,x')
 :=
 \inf_{\gamma:x\leadsto x'}L_{g_O}(\gamma),
\]
with value $+\infty$ if there is no admissible path.

\begin{proposition}[Ambient versus intrinsic image distance]
\label{prop:path}
For all path-connected pairs,
\[
 d_O(x,x')\leq\rho_O(x,x').
\]
Equality holds if a $d_Y$-minimizing geodesic from $O(x)$ to $O(x')$
lies in $O(X)$ and has a lift joining $x$ to $x'$.
\end{proposition}

\begin{proof}
Every image curve joining $O(x)$ to $O(x')$ has length at least their
ambient Riemannian distance. Taking the infimum gives the inequality. A
lift of a minimizing geodesic realizes equality.
\end{proof}

\begin{remark}
The local tensor $g_O$, the path pseudometric $\rho_O$, and the global
pullback pseudometric $d_O=d_Y\circ(O\times O)$ are therefore not
interchangeable without the hypotheses of \cref{prop:path}. This corrects
the ambiguity in v0.1.
\end{remark}

\section{Reference shells and calibrated decisions}

\begin{definition}[Exact and tolerant shells]
Let $r_*>0$ and $\tau\geq0$ have the same scale type as $d_Y$. For
$x\in X$, define
\[
 S_O(x;r_*):=
 \set{x'\in X:d_O(x,x')=r_*},
\]
\[
 S_O^\tau(x;r_*):=
 \set{x'\in X:\abs{d_O(x,x')-r_*}\leq\tau}.
\]
The exact shell is the case $\tau=0$.
\end{definition}

\begin{definition}[Observed distance graph]
For a finite set $A=\{x_1,\ldots,x_n\}\subseteq X$, define
\[
 G_{O,r_*,\tau}(A)
 :=
 \bigl(A,E_{O,r_*,\tau}\bigr),
\]
\[
 \{x_i,x_j\}\in E_{O,r_*,\tau}
 \quad\Longleftrightarrow\quad
 \abs{d_O(x_i,x_j)-r_*}\leq\tau.
\]
\end{definition}

\begin{warning}[Exact equality and measured data]
For continuously noisy data, exact numerical equality is generally not a
stable event. The mathematical shell $\tau=0$ remains valid, but an
empirical classifier must declare a tolerance or a hypothesis test.
\end{warning}

\begin{proposition}[Deterministic error guarantee]
\label{prop:error}
Suppose an estimated distance $\widehat d_{ij}$ obeys
\[
 \abs{\widehat d_{ij}-d_O(x_i,x_j)}\leq u
\]
with a declared bound $u\geq0$. The rule
\[
 \abs{\widehat d_{ij}-r_*}\leq\tau
\]
has the following guarantees:
\begin{enumerate}[label=(\roman*)]
\item every exact-shell pair is accepted if $\tau\geq u$;
\item every accepted pair satisfies
      $\abs{d_O(x_i,x_j)-r_*}\leq\tau+u$.
\end{enumerate}
\end{proposition}

\begin{proof}
Both statements are direct applications of the triangle inequality.
\end{proof}

If $u$ is probabilistic rather than deterministic, the confidence level,
dependence structure, and multiple-testing policy belong to the protocol.

\section{Coherent unit covariance}

Let $a>0$ be a change of numerical length unit. Define
\[
 d_Y'=a\,d_Y,\qquad
 r_*'=a\,r_*,\qquad
 \tau'=a\,\tau,\qquad
 \varepsilon_Y'=a\,\varepsilon_Y.
\]

\begin{theorem}[Unit covariance of shells and graphs]
\label{thm:unit}
\[
 S_{O,d_Y'}^{\tau'}(x;r_*')
 =
 S_{O,d_Y}^{\tau}(x;r_*),
\]
and therefore
\[
 G_{O,r_*',\tau'}(A;d_Y')
 =
 G_{O,r_*,\tau}(A;d_Y).
\]
\end{theorem}

\begin{proof}
\[
\abs{a\,d_O(x,x')-a\,r_*}\leq a\,\tau
\quad\Longleftrightarrow\quad
\abs{d_O(x,x')-r_*}\leq\tau.
\]
\end{proof}

The normalized quantities
\[
 \bar d_O=\frac{d_O}{r_*},
 \qquad
 \bar\tau=\frac{\tau}{r_*},
 \qquad
 \bar\varepsilon_Y=\frac{\varepsilon_Y}{r_*}
\]
are invariant under the same coherent rescaling. A formula
$d_O(x,x')=1$ is admissible only after the convention $r_*=1\,U_Y$ has
been declared. The typed formula is $d_O(x,x')=r_*$.

\begin{definition}[Dimensionless log-scale coordinate]
For a positive reference scale $\ell_*$ and resolution $\varepsilon>0$,
\[
 s_b(\varepsilon;\ell_*):=
 \log_b\frac{\ell_*}{\varepsilon},
 \qquad b>1.
\]
\end{definition}

No logarithm of a bare physical length is admissible.

\section{Matched resolutions under Lipschitz observation}

\begin{definition}[Typed Lipschitz converter]
An observation $O:(X,d_X)\to(Y,d_Y)$ is $L_O$-Lipschitz if
\[
 d_Y(Ox,Ox')\leq L_O\,d_X(x,x')
\]
for all $x,x'$. Semantically,
\[
 [L_O]=U_Y/U_X.
\]
\end{definition}

\begin{proposition}[Resolution propagation]
\label{prop:resolution}
If an input cell has $d_X$-diameter at most $\varepsilon_X$, its image has
$d_Y$-diameter at most $L_O\varepsilon_X$. Hence a conservative matched
output resolution satisfies
\[
 L_O\varepsilon_X\leq\varepsilon_Y.
\]
\end{proposition}

\begin{proof}
Apply the Lipschitz inequality to every pair in the input cell.
\end{proof}

\begin{theorem}[Composition law]
\label{thm:composition}
If $O_1:(X,d_X)\to(Y,d_Y)$ is $L_1$-Lipschitz and
$O_2:(Y,d_Y)\to(Z,d_Z)$ is $L_2$-Lipschitz, then
$O_2\circ O_1$ is $L_2L_1$-Lipschitz. Matched scales propagate as
\[
 \varepsilon_Y\geq L_1\varepsilon_X,\qquad
 \varepsilon_Z\geq L_2\varepsilon_Y
 \quad\Longrightarrow\quad
 \varepsilon_Z\geq L_2L_1\varepsilon_X.
\]
\end{theorem}

\begin{proof}
Substitute the first Lipschitz inequality into the second. The scale
statement follows by transitivity.
\end{proof}

This is the distance-scale bridge to the matched covering defect in the
Metric Entropy v0.2 module.

\section{Orthogonal projection: cylinder and quotient}

Let
\[
 \pi:\R^3\to\R^2,\qquad \pi(x,y,z)=(x,y),
\]
with Euclidean metrics.

\begin{proposition}[Observed shell under projection]
For $X_0=(x_0,y_0,z_0)$ and $r_*>0$,
\[
 S_\pi(X_0;r_*)
 =
 \set{(x,y,z):(x-x_0)^2+(y-y_0)^2=r_*^2}.
\]
Thus the exact observed shell is a circular cylinder of radius $r_*$.
\end{proposition}

\begin{proof}
The condition
$d_\pi(X,X_0)=\norm{\pi(X)-\pi(X_0)}_2=r_*$ is precisely the displayed
equation; $z$ is free.
\end{proof}

\begin{corollary}[Hidden distance profile]
If $d_\pi(X,X')=r_*$, then
\[
 \norm{X-X'}_2^2
 =
 r_*^2+(z-z')^2.
\]
One observed shell radius therefore corresponds to a continuum of hidden
Euclidean distances greater than or equal to $r_*$.
\end{corollary}

The quotient $\R^3/{\sim_\pi}$ is isometric to $\R^2$; the cylinder is the
preimage of a planar circle, not a new type of planar unit-distance
configuration.

\section{Fiber multiplicity and counting rules}

Two non-equivalent counts must be declared:
\begin{enumerate}[label=(\alph*)]
\item \emph{visible count}: replace $A$ by the set of distinct images
      $O(A)$ and count image pairs;
\item \emph{hidden-multiplicity count}: count pairs of distinct hidden
      states whose images satisfy the shell predicate.
\end{enumerate}

\begin{theorem}[Complete bipartite fiber lift]
\label{thm:fibers}
Let $F_a=O^{-1}(a)$ and $F_b=O^{-1}(b)$, with
$d_Y(a,b)=r_*$. If a finite sample contains $m$ states in $F_a$ and $k$
states in $F_b$, its hidden-multiplicity graph contains $K_{m,k}$ and at
least $mk$ shell edges.
\end{theorem}

\begin{proof}
Every cross-fiber pair has observed distance $d_Y(a,b)=r_*$.
\end{proof}

\begin{proposition}[Reduction to the visible problem]
If $O$ is injective on a finite sample $A$ and the count is taken over
distinct images, then
\[
 \abs{E_{O,r_*,0}(A)}
 =
 \#\set{\{a,b\}\subset O(A):d_Y(a,b)=r_*}.
\]
\end{proposition}

\begin{proof}
The restriction $O|_A$ is a vertex bijection preserving the edge
predicate.
\end{proof}

\section{Dimensionally valid periodic and half-step readout}

Let a hidden diameter ladder be
\[
 d_n=d_0+n\Delta d,\qquad n\in\Z,
\]
where $d_0$ and $\Delta d$ are lengths. Let $P>0$ be a physical period and
$\phi_0$ an angle. Define
\[
 h(n)
 =
 \cos\left(
 2\pi\,\frac{d_n-d_0}{P}+\phi_0
 \right).
\]
The trigonometric argument is an angle because $(d_n-d_0)/P$ is
dimensionless.

\begin{theorem}[Rational-step factorization]
\label{thm:alias}
If
\[
 \frac{\Delta d}{P}=\frac pq,\qquad
 p\in\Z,\ q\in\mathbb N,\ \gcd(p,q)=1,
\]
then $h$ factors through $\Z/q\Z$.
\end{theorem}

\begin{proof}
For every $n$,
\[
h(n+q)
=
\cos\left(2\pi\frac{p(n+q)}q+\phi_0\right)
=
\cos\left(2\pi\frac{pn}q+2\pi p+\phi_0\right)
=h(n).
\]
\end{proof}

The theorem states factorization, not injectivity on the $q$ residue
classes; cosine symmetry may identify additional classes.

\begin{corollary}[Half-step parity]
If $\Delta d=P/2$ and $\phi_0=0$, then
\[
 h(n)=\cos(\pi n)=(-1)^n,
\]
so the readout is exactly the parity quotient $\Z\to\Z/2\Z$.
\end{corollary}

This replaces the dimensionally invalid expression
$d_n=d_0+n/2$ unless a length unit for the half-step has first been
declared.

\section{Unit-distance graphs as an application}

For $P\subset\R^2$ finite and $r_*>0$, write
\[
 U_{r_*}(P)
 :=
 \#\set{\{p,q\}\subset P:\norm{p-q}_2=r_*}.
\]
Euclidean similarity $p\mapsto p/r_*$ gives
\[
 \max_{\abs P=n}U_{r_*}(P)
 =
 \max_{\abs P=n}U_1(P).
\]
Thus the classical extremal function is independent of the chosen
positive reference radius after coherent rescaling.

For an observation map $O:X\to\R^2$, define instead
\[
 U_{O,r_*}^{\mathrm{hidden}}(A)
 :=
 \#E_{O,r_*,0}(A).
\]
This quantity can be quadratic because of \cref{thm:fibers}. It is a
different problem: multiplicity inside observation fibers is part of the
data.

\begin{warning}[Erdos firewall]
The present interface proves no new asymptotic upper or lower bound for
the classical planar unit-distance function. It only states the data
needed to define observed-distance variants without dimensional or
counting ambiguity.
\end{warning}

\section{Normative protocol ledger}

For a static theorem with exact metric data, the stochastic-channel field
is recorded as \texttt{explicit-noiseless}, while temporal resolution and
observation horizon are recorded as \texttt{not-applicable}. They become
active fields for time-dependent geometry. This module defines no scalar
defect, so \texttt{defect-normalizations} is recorded as
\texttt{not-applicable}; a downstream distortion or decision loss must
provide its own typed normalization.

\begin{longtable}{L{0.23\textwidth}L{0.18\textwidth}L{0.49\textwidth}}
\toprule
Field & Type & Required declaration\\
\midrule
\endhead
Hidden space & set, manifold, or state space & admissible points and
equivalence before observation\\
Hidden metric $d_X$ & $U_X$ & metric convention and domain\\
Observation $O$ & map $X\to Y$ & deterministic, stochastic, or composed
channel role\\
Noise/channel policy & kernel or exact flag & stochastic law or explicit
noiseless statement\\
Observed metric $d_Y$ & $U_Y$ & ambient or intrinsic convention\\
Observed pseudometric $d_O$ & $U_Y$ & global formula
$d_Y(Ox,Ox')$\\
Reference radius $r_*$ & $U_Y$, positive & shell or incidence target\\
Input resolution $\varepsilon_X$ & $U_X$, positive & quantizer or cover
convention\\
Output resolution $\varepsilon_Y$ & $U_Y$, positive & matched-scale rule\\
Temporal resolution & time or not applicable & required for
time-dependent geometry\\
Observation horizon & time or not applicable & required for a dynamic
trace\\
Tolerance $\tau$ & $U_Y$, nonnegative & exact, bounded-error, or
hypothesis-test policy\\
Uncertainty $u$ or law & $U_Y$ or distribution & confidence/dependence
statement\\
Counting rule & discrete & visible images or hidden multiplicities\\
Reference scale $\ell_*$ & same type as scale variable & normalization and
zero policy\\
Logarithm base $b$ & dimensionless, $b>1$ & every logarithmic trace\\
Estimator & declared map & output, window, and uncertainty summary\\
Defect normalization & typed ratio or not applicable & required before
any scalar loss aggregation\\
\bottomrule
\end{longtable}

\section{Claim audit}

\begin{longtable}{L{0.44\textwidth}L{0.20\textwidth}L{0.26\textwidth}}
\toprule
Statement & Status & Necessary hypotheses\\
\midrule
\endhead
$d_O=d_Y\circ(O\times O)$ is a pseudometric & theorem & $d_Y$ metric,
$O$ a map\\
$X/{\sim_O}$ is isometric to $O(X)$ & theorem & quotient by equal
observations\\
$g_O=O^*g_Y$ determines the same global distance as $d_O$ & not generally
true & requires a suitable lift of minimizing image paths\\
Observed exact shells are unit-covariant & theorem & coherent rescaling
of metric, radius, and tolerance\\
Fiber multiplicity produces $K_{m,k}$ & theorem & two finite sampled
fibers at the target separation\\
Half-step readout recovers the hidden state & false & parity readout is
many-to-one\\
Observation-map graphs solve the classical unit-distance problem & not
claimed & the classical problem is only the injective visible special
case\\
\bottomrule
\end{longtable}

\section{Conclusion}

The common distance-scale object is not a bare number. It is
\[
\bigl(X,d_X,O,Y,d_Y,d_O,r_*,
\varepsilon_X,\varepsilon_Y,\tau,\mathsf U,\text{counting rule}\bigr).
\]
The global pseudometric identifies exactly the observation fibers; its
quotient is the metric image. Physical comparisons use typed radii and
tolerances, logarithms use reference ratios, and resolutions propagate
through declared Lipschitz converters. These rules are the dependency
needed by scale-trace modules such as Mandelbrot rulers.

\begin{thebibliography}{9}

\bibitem{BBI}
D. Burago, Y. Burago, and S. Ivanov,
\emph{A Course in Metric Geometry},
American Mathematical Society, 2001.

\bibitem{Heinonen}
J. Heinonen,
\emph{Lectures on Analysis on Metric Spaces},
Springer, 2001.

\bibitem{Lee}
J. M. Lee,
\emph{Riemannian Manifolds: An Introduction to Curvature},
Springer, 1997.

\bibitem{Erdos}
P. Erdos,
``On sets of distances of $n$ points,''
\emph{American Mathematical Monthly} 53 (1946), 248--250.

\end{thebibliography}

\end{document}
